\section{Source and verification notes}

\subsection{Figures and their data sources}

\begin{itemize}
\item \textbf{Figure 1, pilot instance space:}
  \sourcepath{pilot/work/results/instance_space.csv} and
  \sourcepath{pilot/work/results/perf_matrix.csv}.
\item \textbf{Figure 2, SA/GR ratio by fleet and time window:}
  \sourcepath{pilot/work/results/perf_matrix.csv}.
\item \textbf{Figure 3, generation profile:}
  measured values in \sourcepath{docs/02_tests_and_results.md}, T7. The
  after-fix total wall-time bar is from the measured 115-second rerun. There is
  no source CSV for the profile timings.
\end{itemize}

All three PDFs are regenerated by \sourcepath{make_figures.py}. The seed-count
comparison is intentionally text-only because the available pre-fix data do
not reproduce the documented value.

\subsection{Numbers not independently verified from raw records}

\begin{itemize}
\item The pre-fix count of 6 distinct \texttt{od\_distance\_mean} values is
  reported in the docs but not reproduced by the available data. The available
  pre-fix-labelled CSV has 24 rows and 3 distinct values; the post-fix files
  have 24 distinct values.
\item The rerun exact-model result (5 requests optimal in 2.5 s, 10 requests
  at 25\% gap, and 15 requests not optimal) is reported by session. The source
  does not identify a raw rerun log, nor tabulate rerun objectives or bounds.
\item The 67 MB post-fix maximum-memory result is expressly described as
  suspicious in the source and is not a reliable memory estimate.
\item The final outcome of generation job 29846018, the measured station-matrix
  time and memory, the cause of the extra generation time, and resource sizes
  for subsequent batches are not available.
\item The exact-model objectives and bounds for the rerun, and the comparative
  SA/v2 vehicle-use table, remain unreported.
\end{itemize}

\subsection{Source inconsistencies to resolve}

\begin{itemize}
\item The feature list gives general request-count ranges of 5000--20000 and
  fleet-size ranges of 15--25\%; the current study design instead uses
  100--1000 requests per instance and 5--40\% fleet size. The study-specific
  design is taken from \sourcepath{isa_set/design.json}, but the older feature
  range should be reconciled in the source notes.
\item The paper-settings file describes the pilot period as 6:30--8:30, while
  the generator-change notes describe pilot v2 as having a real 7:30 peak and
  the current design as low-demand starts between 12:00 and 19:30. These may
  refer to different protocol versions; the source set does not fully explain
  the difference.
\item The docs report six pre-fix distinct mean OD distances over 72 pilot
  cases, but available \sourcepath{pilot/work/pilot_v1/results/features.csv}
  contains 24 rows and 3 distinct values. The 72-row files labelled
  \sourcepath{pilot_v2_old} and \sourcepath{results} each contain 24 distinct
  values. The files therefore do not reproduce the stated before-fix/after-fix
  comparison as described.
\item The design JSON names the network as the Capital Region but does not
  encode the Bornholm exclusion in the network string. The decisions and
  network test specify the exclusion, so it is treated as part of the design
  here.
\end{itemize}
